\begin{abstract}
System-strength screening metrics help decide which inverter-based resource (IBR) connections and operating conditions require detailed electromagnetic-transient (EMT) studies. Each metric is specified with its own network model, evaluation point and interaction factor, and the metrics have been compared only on discrete cases.
This paper shows that SCR, WSCR, ESCR and SDSCR are built from the same quantities: the pre-fault voltage, the driving-point impedance and an interaction term.
Unlike earlier comparisons, the metrics are followed through a time-ordered sequence of operating states, and each change is split exactly into these parts, so any difference can be traced to the metrics themselves or to their conventions.
Three IBR plants on a 37-bus benchmark are studied over a week, with load following a public demand profile, intact and under every single line and transformer outage.
Wherever synchronous-generator commitment varies, it sets the weakest hour for all four metrics. Under outages, the metrics follow the change in network impedance; ESCR and SDSCR oppose it only rarely, while ESCR does so far more often when its interaction factor includes generator voltage control.
Screening over time therefore needs a model of commitment, and ESCR values should state how the interaction factor is obtained.
\end{abstract}

\begin{IEEEkeywords}
grid strength, system strength, inverter-based resources, short-circuit ratio, screening, unit commitment, contingency analysis
\end{IEEEkeywords}
